\documentclass[a4paper, 12pt]{article}

\usepackage{utils}

\renewcommand*{\today}{14 November 2025}

\begin{document}

\hotbox{Analyse 3}{CM 13}{\today}

\section{Intégrales généralisées}

Dans ce chapitre, on considère $a \in \R$ et $b \in ]a, +\infty]$.

\begin{definition}
    Soit $f : [a, b[ \to \R$ telle que $f$ est localement intégrable,
    c'est-à-dire que pour tout $c \in [a, b[$, $f$ est intégrable sur $[a, c]$ (autrement dit, $\lim_{c \to b} \int_a^c f(x) \, dx$ existe et est fini).

    On dit que l'intégrale généralisée $\int_a^b f(x) \, dx$ converge et sa limite est appelée \textbf{somme} de l'intégrale.
\end{definition}

\begin{proposition}{}{}
    Soit $f : [a, b[ \to \R$ localement intégrable. Alors pour tout $c \in ]a, b[$
    $$
    \int_a^b f(x) \, dx \text{ converge } \iff \int_c^b f(x) \, dx \text{ converge.}
    $$
    en cas de convergence on a la relation de Chasles généralisée
    $$
    \int_a^b f(x) \, dx = \int_a^c f(x) \, dx + \int_c^b f(x) \, dx.
    $$
\end{proposition}

\begin{demonstration}
    On a pour tout $d \in ]a, b[$
    $$
    \int_a^d f(x) \, dx = \int_a^c f(x) \, dx + \int_c^d f(x) \, dx.
    $$
    et donc $\lim\limits_{d \to b}f(x) \, dx$ existe si et seulement si $\lim\limits_{d \to b} \int_c^d f(x) \, dx$ existe.
    En cas de convergence, on obtient la relation de Chasles en passant à la limite.
\end{demonstration}

\end{document}